% \subsection{Analisis Permasalahan Audit Trail pada DecMed}

Implementasi sistem \textit{Decentralized Medical Record} (DecMed) \cite{DecMed} telah menghasilkan beberapa sumber informasi yang dapat dimanfaatkan sebagai bukti aktivitas (\textit{audit evidence}), yaitu \textit{PatientAccessLog} dan \textit{transaction block} pada jaringan IOTA. Kedua komponen tersebut telah mencatat sebagian aktivitas yang terjadi selama sistem beroperasi. Meskipun demikian, keduanya belum membentuk sebuah sistem \textit{audit trail} yang memadai untuk mendukung proses audit maupun investigasi insiden keamanan.

\textit{PatientAccessLog} digunakan untuk mencatat aktivitas pemberian hak akses dari pasien kepada tenaga medis. Informasi yang tersimpan di dalamnya telah mencakup atribut penting seperti identitas pasien, identitas tenaga medis, waktu pemberian akses (\textit{timestamp}), serta informasi lain yang berkaitan dengan hak akses tersebut. Dari sisi kelengkapan informasi, log ini telah mampu merepresentasikan aktivitas pemberian hak akses dengan baik.

Namun demikian, \textit{PatientAccessLog} masih memiliki beberapa keterbatasan. Pertama, ruang lingkup pencatatan hanya terbatas pada aktivitas pemberian hak akses sehingga berbagai aktivitas penting lainnya, seperti autentikasi, otorisasi, perubahan konfigurasi, maupun interaksi terhadap data rekam medis, belum terdokumentasi. Kedua, \textit{PatientAccessLog} masih dikaitkan secara langsung dengan fungsionalitas pencabutan hak akses (\textit{revoke access}). Ketika hak akses dicabut, informasi pada entri log yang sama akan diperbarui dengan mengubah status akses menjadi \textit{revoked}. Pendekatan tersebut menyebabkan isi log dapat berubah setelah dibuat sehingga tidak memenuhi karakteristik \textit{audit trail} yang seharusnya bersifat \textit{append-only} dan tidak dapat dimodifikasi. Akibatnya, aktivitas pemberian dan pencabutan hak akses tidak terdokumentasi sebagai dua peristiwa yang berdiri sendiri, sehingga kronologi aktivitas menjadi kurang akurat.

Selain \textit{PatientAccessLog}, DecMed juga memanfaatkan \textit{transaction block} pada jaringan IOTA sebagai bukti transaksi yang terjadi di dalam sistem. Setiap pemanggilan fungsi \textit{smart contract} akan menghasilkan sebuah \textit{transaction block} yang tersimpan secara terdistribusi pada jaringan IOTA. Mekanisme ini memberikan jaminan integritas yang tinggi karena data transaksi bersifat \textit{immutable} dan sulit dimodifikasi setelah berhasil dikonfirmasi oleh jaringan.

Meskipun memiliki keunggulan pada aspek integritas, \textit{transaction block} IOTA tidak dirancang sebagai media \textit{audit trail}. Informasi transaksi disimpan dalam struktur data yang berorientasi pada kebutuhan blockchain sehingga kurang mudah dipahami oleh auditor maupun administrator sistem. Selain itu, mekanisme pencarian (\textit{searching}) dan analisis transaksi masih terbatas karena tidak tersedia fasilitas \textit{query} maupun agregasi log yang memadai untuk mendukung proses investigasi insiden. Dari sisi ketersediaan, \textit{transaction block} juga belum dapat menjamin retensi jangka panjang karena data transaksi pada node IOTA dapat dihapus melalui mekanisme \textit{pruning}. Kondisi tersebut menyebabkan bukti digital tidak selalu tersedia ketika dibutuhkan untuk keperluan audit maupun investigasi.

Berdasarkan hasil analisis tersebut, dapat disimpulkan bahwa DecMed telah memiliki beberapa sumber \textit{audit evidence}, namun belum memiliki sebuah sistem \textit{audit trail} yang terintegrasi. Belum tersedia mekanisme yang secara sistematis menentukan aktivitas yang perlu dicatat, menghasilkan \textit{audit record} dari seluruh komponen sistem, serta mengelola audit trail agar mudah dicari, dianalisis, dan dipertahankan integritas maupun ketersediaannya. Oleh karena itu, diperlukan suatu mekanisme yang mampu menghasilkan (\textit{generate}) dan mengelola \textit{audit trail} secara terpusat sehingga dapat menyediakan bukti digital yang andal untuk mendukung proses audit serta investigasi insiden keamanan.